\documentclass[../../main.tex]{subfiles}
\begin{document}
\section{映射度理论}\label{section:ch17-degree-theory}

作为这些思想在拓扑上的一个重要应用，我们现在定义映射度。设 $M,N$ 是两个同为 $n$ 维的、连通的、紧的、定向光滑流形，$F:M\to N$ 是一个光滑映射。

由于 $H^n_{\mathrm{dR}}(N)\cong\mathbb{R}$，我们可以选取一个 $n$-形式 $\omega$ 使得 $\int_N\omega=1$。$F$ 的\emph{映射度}（degree）定义为
\begin{align}
\deg(F)=\int_M F^*\omega.
\end{align}
下面的命题说明这个定义是合理的。

\begin{proposition}[映射度的性质]\label{proposition:ch17-degree-properties}
设 $M,N$ 如上，$F:M\to N$ 光滑。
\begin{enumerate}[(i)]
\item $\deg(F)$ 是一个整数，不依赖于 $\omega$ 的选取。
\item 若 $F$ 与 $G$ 光滑同伦，则 $\deg(F)=\deg(G)$。
\item 若 $F$ 是一个微分同胚，则 $\deg(F)=+1$（保定向）或 $-1$（反定向）。
\item 若 $F:M\to N$，$G:N\to P$，则 $\deg(G\circ F)=\deg(G)\deg(F)$。
\end{enumerate}
\end{proposition}

\begin{figure}[ht]
\centering
\includegraphics[width=0.6\textwidth]{fig-17-4.png}
\caption{映射度}\label{figure:ch17-degree}
\end{figure}

\begin{proof}
因为 $\int_N\omega=1$，若换一个积分等于 $1$ 的形式 $\omega'$，则 $\omega-\omega'$ 在 $N$ 上积分为零，从而是恰当的（因为 $H^{n-1}_{\mathrm{dR}}(N)$ 对紧连通 $N$ 的顶次以下群为零）。于是 $\int_M F^*(\omega-\omega')=0$。整数值可由 Sard 定理得到：取 $N$ 上一个正则值 $q$，在 $q$ 的每个原像处用局部坐标计算，$\deg(F)$ 等于各原像处符号 $+1/-1$ 之和。其余性质由同伦不变性与函子性质得到。
\end{proof}
\end{document}
